%%This is a very basic article template.
%%There is just one section and two subsections.
\documentclass[fromalign=center,fromrule=aftername,fromphone,fromemail,fontsize=10pt,DIV=13]{scrlttr2}
\usepackage[utf8]{inputenc}
\usepackage[T1]{fontenc}
\usepackage[ngerman]{babel}
\usepackage{soul}

\setkomavar{fromname}{Krankenanstalt der Barmherzigkeit -- II. Med. Abteilung}
\setkomavar{fromaddress}{Hartlgasse 16a, A-1240
Wien}
\setkomavar{fromphone}{+43\,555\,328\,745\,0}
\setkomavar{fromemail}{KADBMED2@grissu.org}
\setkomavar{fromlogo}{\begin{tabular}{c}\LARGE\itshape
Grissu.org\\\small\itshape serving the
community\\\small\itshape since
\oldstylenums{1999}\\\end{tabular}}

\setkomavar{place}{Wien}
%\setkomavar{date}{}

\setkomavar{yourref}{--} % Ihr Zeichen
\setkomavar{yourmail}{--} % Ihr Schreiben vom
\setkomavar{myref}{Az 2009/17913-12} % Unser Zeichen
%\setkomavar{customer}{KNr}
%\setkomavar{invoice}{ER-Nr}
\setkomavar{subject}{Betreff: Pat:}

%\setkomavar{}{}


\setkomavar{backaddress}{Krankenanstalt der Barmherzigkeit, Hartlgasse 16a, A-1240
Wien -- AUSTRIA}



                                                               % set fonts
% \renewcommand{\rmdefault}{ugm} % URW Garamond
% %\renewcommand{\sfdefault}{uop} % URW Utopia
% % \renewcommand{\sfdefault}{pag} % Adobe Avantgarde
\usepackage{lmodern}
% \renewcommand{\sfdefault}{phv}
% \renewcommand{\ttdefault}{pcr} 

 \renewcommand{\rmdefault}{lmr}
 \renewcommand{\sfdefault}{phv}
 \renewcommand{\ttdefault}{lmtt}

\setkomafont{fromname}{\fontfamily{phv}\selectfont}
\setkomafont{fromaddress}{\fontfamily{phv}\selectfont\scriptsize}
%\setkomafont{fromphone}{\small}
%\setkomafont{fromemail}{\small}

%%%%%%%%%%%%%%%%%%%%%%%%%%%%%%%%%%%%%%%%%%%%%%%%%%%%%%%%%%%

\begin{document}
\begin{letter}{Name\\
Adresse1\\
Adresse2\\
PLZ}


\opening{Sehr geehrte Damen und Herren!}


Wir berichten über den Aufenthalt unseres Pat. Herrn Peter \caps{Zapfl}, geb. am 13. 12. 1936, welcher vom 30.03.2009 bis 17.04. 2009 an unserer Abteilung lag.

{\bfseries Aus der Anamnese:}

Kindheit: AE, sowie TE, 
WK II: längerer Lazarettaufenthalt wg. Typhus.
Seit 2003 sind ein bis dato nicht insulinpflichtiger Diabetes mell., sowie eine art. Hypertonie bekannt.
2005 Aufenthalt an h.o. Abtlg wg. eines rehirnigen ischämischen zerebralen Insultes, anschl. Rehabaufenthalt in Bd. Schallerbach.
  Zuletzt lag er vom 21.11.2008 bis 23.12.2008 wegen eines nicht insulinpflichtigen, entgleisten Diabetes mell., bei Z. n. Zerebralem ischäm. Insult 2005 mit HSS li und Art. Hypertonie an unserer Abteilung. 
Zwischenzeitlich keine auffälligkeiten.
Alk.: gelegentl.
Nic.: neg.
Coff.: fallweise
Harn, Stuhl: unauff.
Schlaf: unauff.
Ven.: neg
Med.: Renitec 10 mg 1-1-1-
          Diabetex 850 mg 1-1-1, Diab Diät 12 BE 2-1-4-2-2-1
          Thromboass 1-0-0

 Die nunmehrige Aufnahme erfolgte mit dem ASB, der wegen zunehmender Verschlechterung des AZs des Pat. von Wohnungsachbarn berufen wurde.

Auffälliges im Aufnahmestatus: 

 Deutlich herabgesetzter AZ, noch ausreichender EZ, zeitlich und örtlich teilorientiert, motor Unruhe. Akt Beweglichkeit der li OE und UE, sowie Sehnenreflexe li geringer gegenüber re. Haut trocken, Turgor herabgesetzt, kein Decubitus nachweisbar.

Befunde:

Labor siehe Ablichtung.

C/P.: Keine wesentl Änderung gegenüber Vobefund 12.2008 
Schädel CT: Keine wesentliche Änderung gegenüber Letzbefund vom 12.2008, Feinbefund siehe siehe Ablichtung

Augen: Z. f. Erstgradige hypertone Veränderungen der Retinagefässe, kein HW für diabet. Retinopathie.

Neuro: Anamnst Z. n. Insult 2005, mit HSZ li, kein HW auf Reinsult, bzw Spastizität, weiterhin  phys terap. Maßnahmen empf.

Uro: Prost. rect. dig. Unauff.

EKG: SR 74, LT, altersentsprechender Kurvenverlauf.

RR bei mehrfachen Kontrollen gegen Ende des Aufenthaltes normoton, zuvor im hypotonen Bereich.

Epikrise:

Nach Behebung des bei Aufn. Bestehenden Flüssigkeitsdefizites mit geeigneter Infusionstherapie, besserte sich der AZ des Pat. zusehends. Mit in der Folge durchgeführten physikal. Therap.Maßnahmen konnte der Mobilisisationsgrad des Pat. günstig beeinflußt werden. Der Pat. wird mit bis zum selbständigen Gehen mit Rollator mobilisiert entlassen. Die zu Aufenthaltsbeginn durchwegs im hypotonen Bereich gemessenen RR-Werte haben sich in mehrfachen Kontrollen unter dem angegeben antihypertonen Regime normalisiert. Ebeso hat sich die ursprünglich unbefriedigend gefundene diabetische Stoffwechsellage, bei konsequenter Einhaltung des vorgeschlagenen Regimes im therapeutischen Bereich stabilisiert. Im Einverständnis mit dem Pat. wurden zur Gerwährleistung einer regelmäßigen und auch ausreichenden Flüssigkeitszufuhr, sowie einer gesicherten Diabetes Diät, eine Heimhilfe, sowie Essen auf Rädern bestellt. Eine Kontrolle  einer Carotissonographie ist angezeigt. Ein entspr. Termin wurde am 27.12.2009  12,30 Uhr an der Angiolog Ambulanz d. KA Karl Stiftung fixiert. Eine Befundbesprechung ist nach Terminvereinbarung an h.o. Int. Ambulanz möglich.


 Diagnosen:

Hypohydratation 
Diabetes mell 
Art. Hypertonie
Z.n. Zerebralem Insult mit HSS li 2005

Therapie:

Renitec tab. 5,0 mg 1-1-1
Diabetex 850 mg 1-1-1
Diab. Diät 12 BE (2-1-4-2-2-1)
Thromboass 1-0-0
tgl. Flktszufuhr nicht unter 2 l

Ko.:
NFP.  Elektrolyte RR BZ (Pat. führt seit Jahren Selbstkontrollen durch)
weiterhin neurolog. und ophthalmolog. Ko

                                                                                                     Wien,28.09.2009

\closing{Mit freundlichen Grüßen\vspace{2em}\\ Dr. Adolf Schrammel, Stationsarzt}



vidit OA Univ.Doz. Dr. Alois Schremser

\ps
\cc{Hausarzt Dr. Schlögel, Friedrich}
\encl{Labor}

\end{letter}
\end{document}

